\section{Finding colorings and independent sets}

\subsection{Colorable models}

% For a graph on $n$ vertices, we define a $(k, \delta)$-coloring as a list of $k$ disjoint independent sets in $G$ that cover $(1-\delta)n$ of the vertices of $G$.
% Then we define a tractable model as follows.

% \begin{definition}[$(M, \delta)$-coloring]
% Let $M \in \R^{k \times k}_{\geq 0}$ be a reversible row stochastic matrix with zeros on the diagonal.
% For a graph $G(V, E)$, we say that $\chi: V \to [k]$ is a $(M,\delta)$-coloring if $\Norm{M - M(\tilde{A}, \chi)}_{\max} \leq \delta$ and $\chi$ is a $(k, \delta)$-coloring of $G$.
% \end{definition}

% \begin{definition}[Tractable model]
%     \label{def:tract-model}
%     Let $M \in \R^{k \times k}_{\geq 0}$ be a reversible row stochastic matrix with zeros on the diagonal and non-positive second eigenvalue.
%     We say $M$ is tractable if, for every $\e > 0$, there exist $\delta, \lambda > 0$ and a polynomial-time algorithm $\mathcal{A}$ such that, given an $(M, \delta)$-colorable $\lambda$-one-sided expander $G$, the algorithm $\mathcal{A}$ outputs a $(k, \e)$-coloring for $G$.
% \end{definition}

Our most general result for $k$-colorings is the following:
\begin{theorem}[Result for $k$-coloring]
    \label{thm:kalg-partial}
    Let $M \in \R^{k \times k}_{\geq 0}$ be a reversible row stochastic matrix with zero on the diagonal and non-positive second eigenvalue, and let $\pi$ be its stationary distribution.
    Let $S_1, \ldots, S_{k'}$ be a partition of $[k]$ such that, for all $a, b \in S_i$ the rows $M^{(a)}$ and $M^{(b)}$ of $M$ are equal, and for all $a \in S_i, b \in S_j$ with $i \neq j$ the rows $M^{(a)}$ and $M^{(b)}$ of $M$ are distinct.
    For all $i \in [k']$, let $a_i^* = \argmax_{a \in S_i} \pi(a)$, and let $p_i = \max\Paren{0, \pi(a_i^{*}) - \sum_{\substack{a \neq a_i^{*} \in S_i}} \pi(a)}$.
    Then, given a $d$-regular $(M, \delta)$-colorable $\lambda$-one-sided expander $G$, where $\delta, \lambda > 0$ are small enough, there exists an algorithm that runs in time $n^{O(1)} \cdot O_{\delta, \lambda, M}(1)$ and outputs a $(k', 1 - \sum_{i \in [k']} p_i + O_M(\sqrt{\delta} + \lambda + \delta/\lambda))$-coloring for $G$.
\end{theorem}

As a direct corollary, we approach a full coloring in the case when all rows of $M$ are distinct:
\begin{corollary}[Result for $k$-coloring, distinct rows]
    \label{thm:kalg}
    Let $M \in \R^{k \times k}_{\geq 0}$ be a reversible row stochastic matrix with zero on the diagonal, non-positive second eigenvalue, and distinct rows.
    Then, given a $d$-regular $(M, \delta)$-colorable $\lambda$-one-sided expander $G$, where $\delta, \lambda > 0$ are small enough, there exists an algorithm that runs in time $n^{O(1)} \cdot O_{\delta, \lambda, M}(1)$ and outputs a $(k, O_M(\sqrt{\delta} + \lambda + \delta/\lambda))$-coloring for $G$.
\end{corollary}

We start by proving two lemmas that we will use in the proof.
First, we show that the models we consider in \cref{thm:kalg-partial} are well-behaved: they have a unique stationary distribution in which each vertex has non-zero probability.

\begin{lemma}[Model color class size lower bound]
\label{lemma:model-color-size}
Let $M \in \R^{k \times k}_{\geq 0}$ be a reversible row stochastic matrix with zeros on the diagonal.
Then the stationary distribution $\pi$ of $M$ is unique and satisfies $\min_{a \in [k]} \pi_a \geq \Omega_M(1)$.
Furthermore, for every $A \in \R^{n \times n}$ and every $k$-coloring $\chi : [n] \to [k]$ of $[n]$ such that $\Norm{M - M(A, \chi)} \leq \delta$, then also $\Abs{\frac{1}{n}|\chi^{-1}(a)| - \pi_a} \leq 8\sqrt{\delta}k$ for all $a \in [k]$.
\end{lemma}
\begin{proof}
First, because $M$ has non-positive second eigenvalue, its non-zero entries describe a connected graph.
Furthermore, for a stationary distribution $\pi$ of $M$, for all $a \neq b \in [k]$ we have $\pi_a M_{ab} = \pi_b M_{ba}$, so if $M_{ab},M_{ba}>0$ also 
\[\Omega_M(1) \leq \frac{\pi_a}{\pi_b} = \frac{M_{ab}}{M_{ba}} \leq O_M(1)\,.\]
Then, because the non-zero entries of $M$ describe a connected graph, we can chain inequalities of this form to obtain for all $a \neq b \in [k]$ the ratios $\pi_a/\pi_b$.
Together with the fact that $\sum_{a \in [k]} \pi_a = 1$, we get that $\pi$ is uniquely determined by $M$ and satisfies $\min_{a \in [k]} \pi_a \geq \Omega_M(1)$.

Second, we have that $M(A, \chi)_{ab}$ is $\delta$-close to $M_{ab}$ for all $a,b\in [k]$. 
Also, for all $a \neq b \in [k]$ we have $\frac{1}{n}|\chi^{-1}(a)| M(A, \chi)_{ab} = \frac{1}{n}|\chi^{-1}(b)| M(A, \chi)_{ba}$ and $\sum_{a \in [k]} \frac{1}{n}|\chi^{-1}(a)| = 1$.
Then, if $M_{ab},M_{ba}>0$ and $\delta > 0$ is small enough such that $\delta \leq \sqrt{\delta} M_{ab}, \sqrt{\delta} M_{ba}$, also 
\[(1-4\sqrt{\delta}) \frac{\pi_a}{\pi_b} \leq \frac{(1-\sqrt{\delta})M_{ab}}{(1+\sqrt{\delta})M_{ba}} \leq \frac{\frac{1}{n}|\chi^{-1}(a)|}{\frac{1}{n}|\chi^{-1}(b)|} = \frac{M(A, \chi)_{ab}}{M(A, \chi)_{ba}} \leq \frac{(1+\sqrt{\delta})M_{ab}}{(1-\sqrt{\delta})M_{ba}} \leq (1+4\sqrt{\delta})\frac{\pi_a}{\pi_b}\,.\]
This implies that for all $a, b \in[k]$ we have 
\[(1-8\sqrt{\delta}k) \frac{\pi_a}{\pi_b} \leq (1-4\sqrt{\delta})^{k-1} \frac{\pi_a}{\pi_b} \leq \frac{\frac{1}{n}|\chi^{-1}(a)|}{\frac{1}{n}|\chi^{-1}(b)|} \leq (1+4\sqrt{\delta})^{k-1} \frac{\pi_a}{\pi_b} \leq (1+4\sqrt{\delta}k) \frac{\pi_a}{\pi_b}\,,\]
so summing over all $a \in [k]$ we get 
\[(1-8\sqrt{\delta}k)\frac{1}{\pi_b} \leq \frac{1}{\frac{1}{n}|\chi^{-1}(b)|} \leq (1+4\sqrt{\delta}k) \frac{1}{\pi_b}\,,\]
which finally implies that $\Abs{\frac{1}{n}|\chi^{-1}(b)| - \pi_b} \leq 8\sqrt{\delta}k$ for all $b \in [k]$.
\end{proof}

Second, we prove that a $k$-colorable graph that has small variance as per \cref{lemma: expansionimpliesvariance} also has $\tilde{A}$ that is close to $W = Z M(\tilde{A}, \chi) B^{-1} Z^\top$ on the indicators of the color classes.

\begin{lemma}[Small variance implies close indicators]
\label{lem:small-var-to-indic}
Let $G$ be an $n$-vertex graph with adjacency matrix $A$, and let $\tilde{A} = \frac{1}{d}A$ for some $d>0$.
Let $\chi: [n] \to [k]$ be a $k$-coloring of its vertices with $\min_{a \in [k]} |\chi^{-1}(a)| \geq c n$, and consider its partition matrices $Z \in \{0,1\}^{n \times k}$ and $B \in \R^{k \times k}$ as per \cref{def:partition_Z}.
For $x \in [n]$ and $a \in [k]$, let $\D{x}{a} = \sum_{y \in \chi^{-1}(a)} \tilde{A}_{xy}$.
Suppose that, for any $a, b \in [k]$ and $\bm{x}$ uniformly random in $\chi^{-1}(a)$,
\[\E_{\bm{x} \sim \chi^{-1}(a)} \Paren{\D{\bm{x}}{b} - \E_{\bm{x} \sim \chi^{-1}(a)} \D{\bm{x}}{b}}^2 \leq \delta\,.\]
Then, for $W = Z M(\tilde{A}, \chi) B^{-1} Z^\top$ with $M(\tilde{A}, \chi)$ defined as in \cref{def:model} with respect to $\tilde{A}$, it holds for all $a \in [k]$ that
\[\Norm{\tilde{A} \frac{Z_a}{\Norm{Z_a}} - W  \frac{Z_a}{\Norm{Z_a}}}^2 \leq \frac{\delta}{c}\,.\]
\end{lemma}
\begin{proof}
Note that $W_{xy} = \frac{E(\chi^{-1}(\chi(x)), \chi^{-1}(\chi(y)))}{d \Abs{\chi^{-1}(\chi(x))} \Abs{\chi^{-1}(\chi(y))}}$.
Then we get 
\begin{align*}
    \Norm{\tilde{A} Z_a - W Z_a}^2 
    &= \sum_{b \in [k]} \sum_{x \in \chi^{-1}(b)} \Paren{\sum_{y \in \chi^{-1}(a)} \tilde{A}_{x,y} - \frac{1}{d} \frac{E(b,a)}{|\chi^{-1}(b)|} }^2\\
    &= \sum_{b \in [k]} \sum_{x \in \chi^{-1}(b)} \Paren{\D{x}{a} - \E_{\bm{x} \sim \chi^{-1}(b)} \D{\bm{x}}{a}}^2\\
    &= \sum_{b \in [k]} |\chi^{-1}(b)|\; \E_{\bm{x} \sim \chi^{-1}(b)} \Paren{\D{\bm{x}}{a} - \E_{\bm{x} \sim \chi^{-1}(b)} \D{\bm{x}}{a}}^2\\
    &\leq \delta n\,,
\end{align*}
so using that $cn \leq \Norm{Z_a}^2 \leq n$ we get that $\Norm{\tilde{A} \frac{Z_a}{\Norm{Z_a}} - W \frac{Z_a}{\Norm{Z_a}}}^2 \leq \frac{\delta}{c}$.
\end{proof}

Let us now prove \cref{thm:kalg-partial}.

\begin{proof}[Proof of \cref{thm:kalg-partial}]
    We note by \cref{lemma:model-color-size} that $\pi$ is unique and satisfies $\min_{a \in [k]} \pi_a \geq \Omega_M(1)$.

    Denote by $\chi: [n] \to [k]$ the $k$-coloring with respect to which $G$ is $(M, \delta)$-colorable.
    Let $\phi: [k] \to [k']$ map each $a \in [k]$ to the $i \in [k']$ for which $a \in S_i$, and then define $\chi' : [n] \to [k']$ as $\chi'(x) = \phi(\chi(x))$ for all $x \in [n]$.
    That is, $\chi'$ is a $k'$-coloring of the graph that starts with the $k$-coloring $\chi$ and then for each $i \in [k']$ contracts all colors in $S_i$ into a single color.

    We aim to apply \cref{thm:part-recovery} with respect to the $k'$-coloring $\chi'$ and approximately recover the partition of the vertices corresponding to $\chi'$.
    For that, we need to argue that the $k'$-coloring $\chi'$ satisfies the conditions of \cref{thm:part-recovery}.
    Afterward, a simple rounding step will lead to the desired guarantee.

    Let $A$ be the adjacency matrix of the graph, and let $\tilde{A} = \frac{1}{d}A$.
    We define the model matrices $M(\tilde{A}, \chi)$ and $M(\tilde{A}, \chi')$ as in \cref{def:model}.
    We will apply \cref{thm:part-recovery} with matrices $M(\tilde{A}, \chi')$ and $\tilde{A}$ (to play the role of $M$ and $A$ in \cref{thm:part-recovery}, respectively).

    We will use the following two claims about the relation between the models defined with respect to $\chi$ and $\chi'$.
    First, we relate the entries of $M(\tilde{A}, \chi)$ and $M(\tilde{A}, \chi')$.

    \begin{claim}[Closeness of $M(\tilde{A}, \chi)$ and $M(\tilde{A}, \chi')$]
    \label{claim:close-entries-M}
    For any $i, j \in [k']$, we have 
    \begin{equation*}
    M(\tilde{A}, \chi')_{ij} = \frac{\sum_{b \in S_j} |\chi^{-1}(b)|}{|\chi^{-1}(b^*_j)|} M(\tilde{A}, \chi)_{a^*_ib^*_{j}} \pm O_M(\delta)\,.
    \end{equation*}
    \end{claim}
    \begin{proof}
    A calculation shows that, for $i, j \in [k']$,
    \begin{equation}
    \label{eq:model-form}
    M(\tilde{A}, \chi')_{ij} = \sum_{a \in S_i} \frac{|\chi^{-1}(a)|}{\sum_{a' \in S_i} |\chi^{-1}(a')|} \sum_{b \in S_j} M(\tilde{A}, \chi)_{ab}\,.
    \end{equation}
    We note that, for any two rows $M^{(a)}$ and $M^{(a')}$ that are identical, we have that\linebreak $\Norm{M(\tilde{A}, \chi)^{(a)} - M(\tilde{A}, \chi)^{(a')}} \leq O_M(\delta)$.
    Hence, for any $a, a' \in S_i$, we have that $\sum_{b \in S_j} M(\tilde{A}, \chi)_{ab}$ is $O_M(\delta)$-close to $\sum_{b \in S_j} M(\tilde{A}, \chi)_{a'b}$.
    Denote $b_i^* = a_i^*$ for all $i \in [k']$.
    Then, by \cref{eq:model-form} and the above,
    \begin{align*}
    \begin{split}
    \label{eq:mchi}
    M(\tilde{A}, \chi')_{ij}
    &= \sum_{b \in S_j} M(\tilde{A}, \chi)_{a^*_i b} \pm O_M(\delta)\\
    &= \frac{1}{|\chi^{-1}(a^*_i)|} \sum_{b \in S_j} |\chi^{-1}(b)| M(\tilde{A}, \chi)_{b a^*_i} \pm O_M(\delta)\\
    &= \frac{\sum_{b \in S_j} |\chi^{-1}(b)|}{|\chi^{-1}(a^*_i)|} M(\tilde{A}, \chi)_{b^*_{j} a^*_i} \pm O_M(\delta)\\
    &= \frac{\sum_{b \in S_j} |\chi^{-1}(b)|}{|\chi^{-1}(b^*_j)|} M(\tilde{A}, \chi)_{a^*_ib^*_{j}} \pm O_M(\delta)\,,
    \end{split}
    \end{align*}
    where we used that $|\chi^{-1}(a)| M(\tilde{A},\chi)_{ab} = |\chi^{-1}(b)| M(\tilde{A},\chi)_{ba}$ for all $a, b \in [k]$.
    \end{proof}

    Second, let $W = Z M(\tilde{A}, \chi) B^{-1} Z^\top$, with the partition matrices $Z, B$ corresponding to the $k$-coloring $\chi$, and let $W' = Z' M(\tilde{A}, \chi') (B')^{-1} (Z')^\top$, with the partition matrices $Z', B'$ corresponding to the $k'$-coloring $\chi'$.
    Then we also relate the entries of $W$ and $W'$.

    \begin{claim}[Closeness of $W$ and $W'$]
    \label{claim:close-entries-W}
    For any $x, y \in [n]$, we have 
    \begin{equation*}
    W'_{xy} = W_{xy} \pm O_M(\delta/n)\,.
    \end{equation*}
    \end{claim}
    \begin{proof}
    A calculation shows that, for $x,y \in [n]$,
    \[W_{xy} = \frac{M(\tilde{A}, \chi)_{\chi(x)\chi(y)}}{|\chi^{-1}(\chi(y))|}\] 
    and, by \cref{claim:close-entries-M}, 
    \begin{align*}
    W'_{xy}
    &= \frac{M(\tilde{A}, \chi')_{\chi'(x)\chi'(y)}}{|(\chi')^{-1}(\chi'(y))|}\\
    &= \frac{ \frac{\sum_{b \in S_{\chi'(y)}} |\chi^{-1}(b)|}{|\chi^{-1}(b_{\chi'(y)})|} M(\tilde{A}, \chi)_{a_{\chi'(x)}b_{\chi'(y)}}}{\sum_{b \in S_{\chi'(y)}} |\chi^{-1}(b)|} \pm O_M(\delta/n)\\
    &= \frac{ M(\tilde{A}, \chi)_{a_{\chi'(x)}b_{\chi'(y)}}}{|\chi^{-1}(b_{\chi'(y)})|} \pm O_M(\delta/n)\,.
    \end{align*}
    In \cref{claim:close-entries-M} $a^*_{\chi'(x)}$ and $b^*_{\chi'(y)}$ could in fact be any arbitrary colors in $S_{\chi'(x)}$ and $S_{\chi'(y)}$, respectively.
    Then, letting $a_{\chi'(x)} = \chi(x)$ and $b_{\chi'(y)} = \chi(y)$, we get that $W'_{xy} = W_{xy} \pm O_M(\delta/n)$.
    \end{proof}

    Let us prove now that we satisfy the conditions of \cref{thm:part-recovery} with matrices $M(\tilde{A}, \chi')$ and $\tilde{A}$.
    First, we have trivially that $(B')^{1/2} M(\tilde{A}, \chi') (B')^{-1/2}$ is symmetric and that $\Norm{M(\tilde{A}, \chi')} = 1$.
    We also have that $|\chi^{-1}(a)| \geq \Omega_M(n)$ for all $a \in [k]$ by \cref{lemma:model-color-size}.
    
    \paragraph{Row separation of $M(\tilde{A}, \chi')$}
    We prove now that the rows of $M(\tilde{A}, \chi')$ are separated.
    Consider two rows $M(\tilde{A}, \chi')^{(i)}$ and $M(\tilde{A}, \chi')^{(i')}$ with $i \neq i' \in [k']$.
    Because $M^{(a^*_i)} \neq M^{(a^*_{i'})}$, we have that $\Norm{M^{(a^*_i)} - M^{(a^*_{i'})}} \geq \Omega_M(1)$, so then by the closeness of $M$ and $M(\tilde{A}, \chi)$ we also have that $\Norm{M(\tilde{A}, \chi)^{(a^*_i)} - M(\tilde{A}, \chi)^{(a^*_{i'})}} \geq \Omega_M(1)$.
    Therefore, there exists some $b \in [k]$ such that\linebreak $\Abs{M(\tilde{A}, \chi)_{a^*_ib} - M(\tilde{A}, \chi)_{a^*_{i'}b}} \geq \Omega_M(1)$.
    By \cref{claim:close-entries-M} this implies that, for $j = \phi(b)$, we also have $\Abs{M(\tilde{A}, \chi')_{i j} - M(\tilde{A}, \chi')_{i' j}} \geq \Omega_M(1)$.
    Hence, $\Norm{M(\tilde{A}, \chi')^{(i)} - M(\tilde{A}, \chi')^{(i')}} \geq \Omega_M(1)$.

    \paragraph{Closeness of $\tilde{A}$ and $W'$}
    By definition, we only consider $G$ and colorings with respect to which it has a vertex cover of all monochromatic edges of size at most $\delta n$, so the number of monochromatic edges is bounded by $\delta d n$.
    By \cref{lemma:model-color-size}, the size of each color class is lower bounded by some $\Omega_M(n)$.
    Then by \cref{lemma: expansionimpliesvariance} and \cref{lem:small-var-to-indic} we have that $\Norm{\tilde{A} \frac{Z_a}{\Norm{Z_a}} - W \frac{Z_a}{\Norm{Z_a}}}^2 \leq O_M\Paren{\lambda + \frac{\delta}{\lambda}}$ for all $a \in [k]$.

    By \cref{claim:close-entries-W}, we have that $\Norm{W'-W} \leq \Norm{W'-W}_F \leq O_M(\delta)$.
    Then, for each $S_i$, denoting its elements as $S_i = \{a_1, \ldots, a_{\ell}\}$, we are interested in the closeness of $\tilde{A}$ and $W'$ with respect to the indicator $\sum_{j=1}^{\ell} Z_{a_j}$.
    We have
    \[\Norm{(\tilde{A} - W') \sum_{j=1}^{\ell} Z_{a_j}} \leq \sum_{j=1}^{\ell} \Norm{(\tilde{A} - W') Z_{a_j}} \leq \sum_{j=1}^{\ell} \Norm{(\tilde{A} - W) Z_{a_j}} + \sum_{j=1}^{\ell}\Norm{(W'-W)Z_{a_j}}\,.\]
    For the first term on the right-hand side, we have that $\Norm{\tilde{A} \frac{Z_a}{\Norm{Z_a}} - W \frac{Z_a}{\Norm{Z_a}}}^2 \leq O_M\Paren{\lambda + \frac{\delta}{\lambda}}$ for all $a \in [k]$, so using that $\Norm{Z_a} \leq \sqrt{n}$ we can bound
    \[\sum_{j=1}^{\ell} \Norm{(\tilde{A} - W) Z_{a_j}} \leq \sqrt{n} \cdot O_M\Paren{\sqrt{\lambda + \frac{\delta}{\lambda}}}\,.\]
    For the second term, we have that
    \[\sum_{j=1}^{\ell} \Norm{(W'-W)Z_{a_j}} \leq \sum_{j=1}^{\ell} \Norm{W'-W} \Norm{Z_{a_j}} \leq \sqrt{n} \cdot O_M(\delta)\,.\]
    Then, also using that $\Norm{\sum_{j=1}^{\ell} Z_{a_j}} \geq \Omega_M(\sqrt{n})$, we have
    \[\Norm{(\tilde{A} - W') \frac{\sum_{j=1}^{\ell} Z_{a_j}}{\Norm{\sum_{j=1}^{\ell} Z_{a_j}}}}^2 \leq O_M \Paren{\lambda + \frac{\delta}{\lambda}}\,,\]
    so we satisfy the second condition in \cref{thm:part-recovery} with constant $O_M \Paren{\lambda + \frac{\delta}{\lambda}}$.

    \paragraph{Low rank of $\tilde{A}$}
    Because $\rank_{\geq \lambda}(\tilde{A}) \leq 1$, by applying \cref{cor:spectral_small_to_small} with $\sigma=\lambda$, we get that the bottom threshold rank of $\tilde{A}$ is bounded by $\rank_{\leq -\sqrt{2\lambda}}(\tilde{A}) \leq 1/\lambda^2$.
    Note that, using that $\lambda \leq 1$, we have that $\lambda \leq \sqrt{2 \lambda}$.
    Then, for the purposes of \cref{thm:part-recovery}, we apply the theorem with eigenvalue threshold $\lambda' = \max\{\sqrt{2\lambda}, \Omega_M(1)\}$ for some small enough $\Omega_M(1)$ such that the conditions of the theorem are satisfied, and this allows us to take $t = O(1/(\lambda')^2) = O_M(1)$.

    \paragraph{Rounding}
    Then, by \cref{thm:part-recovery} applied with matrices $M(\tilde{A}, \chi')$ and $\tilde{A}$, we can recover in time $n^{O(1)} \cdot O_{\delta,\lambda,M}(1)$ a list of $O_{\delta,\lambda,M}(1)$ $k'$-colorings $\hat{\chi'}: [n] \to [k']$ such that, for at least one of them, it holds up to a permutation of the color classes that $\mathbb{E}_{\bm{x} \sim \operatorname{Unif}([n])} \1[\hat\chi'(x) \neq \chi'(x)] \leq O_M(\lambda+\delta/\lambda)$.
    Let $\e = O_M(\lambda + \delta/\lambda)$ be equal to this upper bound.
    At this point, for each candidate $k'$-coloring, we apply the rounding algorithm in \cref{lemma: generalrounding} on each of its color classes to obtain $k'$ independent sets, and return the set of $k'$ independent sets that cover most vertices of $G$.

    Let us analyze the approximation guarantee.
    For the best $k'$-coloring, the rounding algorithm in \cref{lemma: generalrounding} will remove at most $2\e n$ vertices due to misclassifications and at most $2\delta n$ vertices due to monochromatic edges with respect to the $k$-coloring $\chi$.
    The rounding algorithm may also need to remove a significant portion of vertices because we recover $k'$-colorings instead of $k$-colorings --- this is necessary even if the $k'$-coloring $\chi'$ is perfectly recovered.
    To analyze these removals, let us start by defining $\bar{a}_i^* = \argmax_{a \in S_i} |\chi^{-1}(a)|$ for each $i \in [k']$.
    Then, for each $i \in [k']$, the rounding algorithm may remove at most $\max\Paren{\sum_{a \in S_i} |\chi^{-1}(a)|, 2\sum_{\substack{a \neq a_i^{*} \in S_i}} |\chi^{-1}(a)|}$ vertices from $(\chi')^{-1}(i)$.
    Let us define for each $i \in [k']$ the minimum fraction of vertices not removed as $\bar{p}_i = \frac{1}{n} \max\Paren{0, |\chi^{-1}(a_i^*)| - \sum_{\substack{a \neq a_i^{*} \in S_i}} |\chi^{-1}(a)|}$.
    Then the algorithm returns a $(k', 1 - \sum_{i \in [k']} \bar{p}_i + 2\e)$-coloring for $G$.
    Finally, by \cref{lemma:model-color-size}, for $\delta > 0$ small enough, $|\chi^{-1}(a)|/n$ and $\pi_a$ are $O(\sqrt{\delta}k)$-close for all $a \in [k]$, so the algorithm also returns a $(k', 1 - \sum_{i \in [k']} p_i + O_M(\sqrt{\delta}) + 2\e)$-coloring for $G$, where $p_i$ is defined as in the statement of the theorem.
    Recalling that $\e = O_M(\lambda + \delta/\lambda)$ gives the stated result.
\end{proof}

\subsection{$2$-colorable models}


Assuming the UGC conjecture, \cite{MR2648426} showed that it is NP-hard to recover even a small linear-sized independent set from a graph that is almost bipartite. As a $k=2$ corollary of \cref{thm:kalg}, we show that assuming expansion additionally, it is possible to receover almost the full bipartition.

\andor{TODO: note that only reversible row stocastic 2 by 2 M is 0 1 1 0}


\andor{TODO: row separation is $1$}

\begin{lemma}[Balanced bipartition]
    Let $G$ be an $n$-vertex $d$-regular graph with adjacency matrix $A$, let $\tilde{A} = \frac{1}{d}A$, and let $M = M(\tilde{A}, \chi)$ be defined as in \cref{def:model} with respect to $\tilde{A}$.
    Let $\chi: [n] \to [2]$ be a $2$-coloring of its vertices, and suppose that $G$ has a vertex cover of all monochromatic edges of size at most $\delta n$.
    Then, $|\chi^{-1}(0)|,\ |\chi^{-1}(0)|  \geq (1/2-\delta)n$, and $M_{12},\ M_{21} \geq 1-4\delta$.
\end{lemma}


\begin{theorem}[Result for $2$-coloring]
\label{thm:2alg}
Let $G$ be an $n$-vertex $d$-regular graph with adjacency matrix $A$ and let $\tilde{A} = \frac{1}{d}A$.
Let $\lambda$ be the second largest eigenvalue of $\tilde{A}$.
Let $\chi: [n] \to [2]$ be a $2$-coloring of its vertices, and suppose that $G$ has a vertex cover of all monochromatic edges of size at most $\delta n$.
Then, for $\delta, \gamma, \lambda > 0$ small enough, there exists an algorithm that runs in time $n^{O(1)} \cdot O_{\delta,\lambda}(1)$ and outputs a $(3, O_{\delta,\lambda}(1))$-coloring for $G$.
\end{theorem}
\begin{proof}
For simplicity of notation, let $M = M(\tilde{A}, \chi)$ be defined as in \cref{def:model} with respect to $\tilde{A}$.
We aim again to apply \cref{thm:part-recovery} with matrices $M$ and $\tilde{A}$, and we verify that its conditions hold.
We have by \cref{lemma: 3coloringrowseparation} that any two distinct rows of $M$ have separation $\gamma$.
Furthermore, the bound $\max_{a \in [3]} \frac{1}{n} |\chi^{-1}(a)| \leq 1/2 - \gamma$ implies that $|\chi^{-1}(a)| \geq 2\gamma n$ for all $a \in [3]$.
For the rest of the proof, we will largely repeat the analysis in the proof of \cref{thm:kalg-partial}, but using that any two distinct rows of $M$ are separated and taking into account explicitly the parameter $\gamma$.

Let $W = Z M B^{-1} Z^\top$, with the partition matrices $Z, B$ corresponding to the $k$-coloring $\chi$.
We have trivially that $B^{1/2} M B^{-1/2}$ is symmetric and that $\Norm{M} = 1$.
We also have that $|\chi^{-1}(a)| \geq 2\gamma n$ for all $a \in [3]$.

In addition, as argued above, we have that any two distinct rows of $M$ have separation $\gamma$.
For the closeness of $\tilde{A}$ and $W$, we have by \cref{lemma: expansionimpliesvariance} and \cref{lem:small-var-to-indic} that $\Norm{\tilde{A} \frac{Z_a}{\Norm{Z_a}} - W \frac{Z_a}{\Norm{Z_a}}}^2 \leq O\Paren{\lambda/\gamma + \frac{\delta}{\gamma \lambda}}$ for all $a \in [3]$.
For the low rank of $\tilde{A}$, we have by applying \cref{cor:spectral_small_to_small} with $\sigma=\lambda$ that $\rank_{\leq -\sqrt{2\lambda}}(\tilde{A}) \leq 1/\lambda^2$, and also $\lambda \leq \sqrt{2 \lambda}$. 

Then we apply \cref{thm:part-recovery} with matrices $M$ and $\tilde{A}$ and eigenvalue threshold\linebreak $\lambda' = \max\{\sqrt{2 \lambda}, \Omega(\gamma^{3/2})\}$, such that it satisfies the condition of the theorem, and such that we can take $t = O(1/(\lambda')^2) = O(1/\gamma^3) = O_{\gamma}(1)$.
Then, we can recover in time $n^{O(1)} \cdot O_{\delta,\gamma,\lambda}(1)$ a list of $O_{\delta,\gamma,\lambda}(1)$ $3$-colorings $\hat{\chi}: [n] \to [3]$ such that, for at least one of them, it holds up to a permutation of the color classes that $\mathbb{E}_{\bm{x} \sim \operatorname{Unif}([n])} \1[\hat\chi(x) \neq \chi(x)] \leq O(\lambda/\gamma^7+\delta/(\lambda\gamma^7))$.
As in the proof of \cref{thm:kalg-partial}, for each candidate $3$-coloring we apply the rounding algorithm in \cref{lemma: generalrounding} on each of its color classes to obtain $3$ independent sets, and return the set of $3$ independent sets that cover most vertices of $G$.
The rest of the analysis is identical to that in the proof of \cref{thm:kalg-partial}.
\end{proof}


\andor{TODO: story about how almost bipartition is actually hard}

\begin{theorem}[Result for $2$-coloring]
    \label{thm:2alg}
    \andor{TODO}
\end{theorem}

\subsection{$3$-colorable models}

For the case $k=3$, it turns out that the model has non-positive eigenvalue and distinct rows whenever its stationary distribution is strictly bounded pointwise by $1/2$. 

The first observation is that we can compute explicitly the entries of $M$ in terms of the stationary distribution.

\begin{lemma}[Model entries for $3$-colorings]
\label[lemma]{lemma: 3coloringfixedaveragedegree}
Let $M \in \R^{3 \times 3}_{\geq 0}$ be a reversible row stochastic matrix with zeros on the diagonal and stationary distribution $\pi$.
Then for all $a \neq b \in [3]$ 
\[M_{ab} = \frac{(2\pi(a)+2\pi(b))-1}{2\pi(a)}.\]
\end{lemma}
\begin{proof}
    For all $a \in [3]$, we have $\sum_{b\in [3]}M_{ab}=1$.
    Also for all $b\in [3]\setminus\Set{a}$, we have $\pi(a) M_{ab} = \pi(b) M_{ba}$.
    Hence, taking the off-diagonal entries of the $M$ matrix as unknown, we have $6$ variables and $6$ independent linear equations, so there is a unique solution.
    Using $\sum_{a\in[3]}\pi(a)=1$, for $a\neq b\neq c\in[3]$ we can express it as
    $$M_{ab}= \frac{(2\pi(a)+2\pi(b))-1}{2\pi(a)}\,.$$
\end{proof}

We then use this observation to deduce that, whenever the largest probability of a color class is upper bounded by $1/2-\gamma$, any two distinct rows of $M$ have separation $\gamma$.

\begin{lemma}[Row separation for $3$-colorings]
\label[lemma]{lemma: 3coloringrowseparation}
Let $M \in \R^{3 \times 3}_{\geq 0}$ be a reversible row stochastic matrix with zero on the diagonal and stationary distribution $\pi$ that satisfies $\max_{a \in [3]} \pi_a \leq 1/2 - \gamma$.
Then any two distinct rows $M^{(a)}$ and $M^{(b)}$ of $M$ satisfy $\Norm{M^{(a)} - M^{(b)}} \geq \gamma$.
\end{lemma}
\begin{proof}
It suffices to prove that all non-diagonal entries of $M$ are bounded away from $0$.
From \cref{lemma: 3coloringfixedaveragedegree}, we get for $a \neq b$ that $M_{ab} = \frac{(2\pi_a+2\pi_b)-1}{2\pi_a}$.
Because $\max_{a \in [3]} \pi_a \leq 1/2-\gamma$, it follows that $\pi_a+\pi_b \geq 1/2+\gamma$, so $M_{ab} \geq \frac{\gamma}{\pi_a} \geq \gamma$.
Then the conclusion follows because $\Norm{M^{(a)} - M^{(b)}} \geq M_{ab} - M_{bb} = M_{ab} \geq \gamma$.
\end{proof}

Finally, we state and prove our result specialized to $3$-colorings.

\begin{theorem}[Result for $3$-coloring]
\label{thm:3alg}
Let $G$ be an $n$-vertex $d$-regular graph with adjacency matrix $A$, and let $\tilde{A} = \frac{1}{d}A$.
Let $\lambda$ be the second largest eigenvalue of $\tilde{A}$.
Let $\chi: [n] \to [3]$ be a $3$-coloring of its vertices, and suppose that $G$ has a vertex cover of all monochromatic edges of size at most $\delta n$.
Suppose $\max_{a \in [3]} \frac{1}{n} |\chi^{-1}(a)| \leq 1/2 - \gamma$.
Then, for $\delta, \gamma, \lambda > 0$ small enough, there exists an algorithm that runs in time $n^{O(1)} \cdot O_{\delta,\gamma,\lambda}(1)$ and outputs a $(3, O(\lambda/\gamma^7+\delta/(\lambda\gamma^7)))$-coloring for $G$.
% A reversible row stochastic matrix $M \in \R^{3 \times 3}_{\geq 0}$ with zero on the diagonal is tractable if its stationary distribution $\pi$ satisfies $\max_a \pi_a < 1/2$.
\end{theorem}
\begin{proof}
For simplicity of notation, let $M = M(\tilde{A}, \chi)$ be defined as in \cref{def:model} with respect to $\tilde{A}$.
We aim again to apply \cref{thm:part-recovery} with matrices $M$ and $\tilde{A}$, and we verify that its conditions hold.
We have by \cref{lemma: 3coloringrowseparation} that any two distinct rows of $M$ have separation $\gamma$.
Furthermore, the bound $\max_{a \in [3]} \frac{1}{n} |\chi^{-1}(a)| \leq 1/2 - \gamma$ implies that $|\chi^{-1}(a)| \geq 2\gamma n$ for all $a \in [3]$.
For the rest of the proof, we will largely repeat the analysis in the proof of \cref{thm:kalg-partial}, but using that any two distinct rows of $M$ are separated and taking into account explicitly the parameter $\gamma$.

Let $W = Z M B^{-1} Z^\top$, with the partition matrices $Z, B$ corresponding to the $k$-coloring $\chi$.
We have trivially that $B^{1/2} M B^{-1/2}$ is symmetric and that $\Norm{M} = 1$.
We also have that $|\chi^{-1}(a)| \geq 2\gamma n$ for all $a \in [3]$.

In addition, as argued above, we have that any two distinct rows of $M$ have separation $\gamma$.
For the closeness of $\tilde{A}$ and $W$, we have by \cref{lemma: expansionimpliesvariance} and \cref{lem:small-var-to-indic} that $\Norm{\tilde{A} \frac{Z_a}{\Norm{Z_a}} - W \frac{Z_a}{\Norm{Z_a}}}^2 \leq O\Paren{\lambda/\gamma + \frac{\delta}{\gamma \lambda}}$ for all $a \in [3]$.
For the low rank of $\tilde{A}$, we have by applying \cref{cor:spectral_small_to_small} with $\sigma=\lambda$ that $\rank_{\leq -\sqrt{2\lambda}}(\tilde{A}) \leq 1/\lambda^2$, and also $\lambda \leq \sqrt{2 \lambda}$. 

Then we apply \cref{thm:part-recovery} with matrices $M$ and $\tilde{A}$ and eigenvalue threshold\linebreak $\lambda' = \max\{\sqrt{2 \lambda}, \Omega(\gamma^{3/2})\}$, such that it satisfies the condition of the theorem, and such that we can take $t = O(1/(\lambda')^2) = O(1/\gamma^3) = O_{\gamma}(1)$.
Then, we can recover in time $n^{O(1)} \cdot O_{\delta,\gamma,\lambda}(1)$ a list of $O_{\delta,\gamma,\lambda}(1)$ $3$-colorings $\hat{\chi}: [n] \to [3]$ such that, for at least one of them, it holds up to a permutation of the color classes that $\mathbb{E}_{\bm{x} \sim \operatorname{Unif}([n])} \1[\hat\chi(x) \neq \chi(x)] \leq O(\lambda/\gamma^7+\delta/(\lambda\gamma^7))$.
As in the proof of \cref{thm:kalg-partial}, for each candidate $3$-coloring we apply the rounding algorithm in \cref{lemma: generalrounding} on each of its color classes to obtain $3$ independent sets, and return the set of $3$ independent sets that cover most vertices of $G$.
The rest of the analysis is identical to that in the proof of \cref{thm:kalg-partial}.
\end{proof}

We also provide a completely different view on the $3$-coloring case, highlighting a potentially deeper connection. Theorem 1.1 of \cite{barak-raghavendra-steurer} shows a general result: given a satisfiable CSP with constraints over the edges of a graph with low \textit{top} threshold rank, a semidefinite programming hierarchy rounding algorithm can efficiently find an assignment to the vertices safisfying most of the constraints. As the colouring constraint is a special case of this, the result implies that a $k$-colouring violating only a small fraction of edges can be recovered efficiently for any $k$ -- what is more, the algorithm doesn't use the bottom eigenspace.

However, we are after small \text{vertex} violation, not small edge violation: we want that the number of vertices that are not colored correctly is small. For $k=3$ we show a surprising connection: under one-sided expansion and balancedness, low edge-violation implies low vertex violation -- hence, an alternative of \cref{thm:3alg} immediately follows from Theorem 1.1 of \cite{barak-raghavendra-steurer}. For convenience, we show the contrapositive instead: high vertex violation (i.e. low agreement) implies high edge violation.



\begin{theorem}[Vertex vs edge violation]\label{thm: vertexvsedgeviolation}
    Let $0<\gamma, \delta<\frac{1}{6}$ be constants. Let $G$ be an $n$-vertex $d$-regular graph with $\lambda_2\leq \frac{\delta^6\gamma^9}{10^2}$ that admits a \tcl $\chi$ such that $\max_{a \in [3]} \frac{1}{n} |\chi^{-1}(a)| \leq 1/2 - \gamma$. Let $\chi'$ be an alternative (not necessarily proper) \tcl with $\text{agree}(\chi, \chi')\leq 1-\delta$. Then, $\chi'$ must have at least $\frac{1}{10^4}\delta^3\gamma^4 nd$ monochromatic edges. 
\end{theorem}

The proof is quite calculation heavy and the calculations are similar to those in \cref{lemma: expansionimpliesvariance} so we only provide a sketch highlighting the main ideas.


\begin{proof}[Proof sketch of \cref{thm: vertexvsedgeviolation}]
    Let $\psi^{-1}(a, b) := \chi^{-1}(a) \cap \chi'^{-1}(b)$ be the intersection of the colour classes of $\chi$ and $\chi'$ -- without loss of generality, we can assume that the identity permutation maximises agreement. 

    The proof starts with a technical lemma showing that as $\chi$ and $\chi'$ have low agreement, there exist colours $a$ and $b$ such that $\Abs{\chi'^{-1}(b)}$, $\Abs{\psi^{-1}(a, b)}$, and $\Abs{\psi^{-1}(b, b)}$ are not too large and not too small either.

    By \cref{lemma: expansionimpliesvariance}, we have that variance of the degrees between $\chi^{-1}(a)$ and $\chi^{-1}(b)$ is small. On the other hand, note that if the average degree between $\psi^{-1}(a, b)$ and $\chi'^{-1}(b)$ was much less than the average degree between $\chi^{-1}(a)$ and $\chi^{-1}(b)$, it would imply that the variance of the degrees between $\chi^{-1}(a)$ and $\chi^{-1}(b)$ is large.
    
    However, if $\chi'$ had only few monochromatic edges, then the average degree between $\psi^{-1}(a, b)$ and $\psi^{-1}(b, b)$ must be small, so the only way the average degree between $\psi^{-1}(a, b)$ and $\chi'^{-1}(b)$ can be large enough is if much more edges from $\psi^{-1}(a, b)$ go to $\chi'^{-1}(b) \setminus \psi^{-1}(b, b)$ than to $\psi^{-1}(b, b)$.

    In this case, however, for suitably chosen $\alpha$ and $\beta$, we can construct a test vector $v$ (orthogonal to $\1$) showcasing that expansion must be poor by lower bounding $\frac{v^\top A_G v}{v^\top v}$ and using the Courant-Fischer Theorem:

    \begin{itemize}
        \item for $x \in \psi^{-1}(a, b)$, let $v_x := 1$
        \item for $x \in \psi^{-1}(b, b)$, let $v_x := -\alpha$
        \item for $x \in \chi'^{-1}(b) \setminus \psi^{-1}(b, b)$, let $v_x := \beta$
        \item otherwise, let $v_x := 0$
    \end{itemize}

    This is a contradiction, so $\chi'$ must have many monochromatic edges.
\end{proof}


\subsection{Random planting}

Instead of a regular $k$-colorable one-sided expander, in this section we consider randomly planting a \kcl in a regular one-sided expander.
To aid presentation, let us first define the model and some of the notation that we will use throughout this section.

\begin{definition}[Randomly planted \kcl] \label{def: randomplanted}
Let $H$ be an $n$-vertex $d$-regular graph. 
Then we say the graph $G$ is obtained by randomly planting a \kcl in $H$ if it is sampled according to the following procedure:
\begin{enumerate}
    \item Sample $\chi: [n] \to [k]$ by sampling $\chi(x) \sim \operatorname{Unif}([k])$ independently for each $x \in [n]$,
    \item Let $G$ be a copy of $H$, and then remove from $G$ each edge $\{u,v\}$ for which $\chi(u) = \chi(v)$.
\end{enumerate}
\end{definition}

% \rares{What to do with these?}
% Letting $d':=\frac{k-1}{k}d$, we define a normalized adjacency matrix $\tilde{A} := \frac1{d'}A_G$ of $G$. For $x \in [n]$ and $a \in [k]$, let $\D{x}{a} = \sum_{y \in \chi^{-1}(a)} \tilde{A}_{xy}$. Then, we define the model $M \in \R^{k \times k}$ such that for any $a, b \in [k]$, has entries $M_{ab} = \frac{1}{|\chi^{-1}(a)|} \sum_{x \in \chi^{-1}(a)}\D{x}{b}$\andor{is it fine to use $M$ instead of $M(\tilde{A},\chi)$?}.
% Let $\rank_{\le x}(H):=\rank_{\le x}\Paren{\frac1{d}A_H}$, and $\rank_{\ge x}(H):=\rank_{\ge x}\Paren{\frac1{d}A_H}$.\andor{is using $\rank{\tilde{A}}$ but $\rank(H)$ fine?}

The following is our main recovery result.
We note that in \cref{thm: random-planting-full}, for the case when the host graph is a one-sided expander, we give an algorithm that recovers a coloring on \emph{all} vertices of the graph.

\begin{theorem}[Result for random planting, partial recovery]
    \label{thm: random-planting-partial}
    Let $H$ be an $n$-vertex $d$-regular graph with adjacency matrix $A_H$.
    Suppose $\rank_{\geq \lambda}(\frac{1}{d}A_H) \leq t$ for some $0 < \lambda < 1$ and $t \in \mathbb{N}$.
    Let $G$ be the graph obtained by randomly planting a \kcl in $H$, and let $\chi: [n] \to [k]$ be the associated \kcl.
    Then, for $\lambda > 0$ small enough, there exists an algorithm that runs in time $n^{O(1)} \cdot \Paren{O_{k}(1)}^{t}$ and outputs with high probability a $(k, O_k(1/d^{\Omega(1)}))$-coloring for $G$.
    %  a list of $\Paren{O_{k}(1)}^{t}$ $k$-colorings $\hat{\chi}: [n] \to [k]$ such that with high probability, for at least one of them, it holds up to a permutation of the color classes that $\mathbb{E}_{\bm{x} \sim \operatorname{Unif}([n])} \1[\hat\chi(\bm{x}) \neq \chi(\bm{x})] \leq O_k(1/d^{\Omega(1)})$.
    % Let $\Omega(1)\leq\lambda\leq \frac1{O(k^3)}$\andor{this $\Omega(1)$ assumption should ideally be removed}, $t\geq 1$, $d\geq 1$ and $k=O(1)$ be parameters. Let $H$ be an $n$-vertex $d$-regular graph with $\rank_{\ge \lambda}(H)\leq t$, and let $G$ be the graph obtained from $H$ by planting a \kcl in it at random.
\end{theorem}

To prove \cref{thm: random-planting-partial}, we first prove some properties of the graph $G$ obtained after randomly planting a \kcl in $H$.

\begin{lemma}[Color class concentration in random planting]
    \label[lemma]{fact: balancedparts}
    Let $H$ be an $n$-vertex graph.
    Let $G$ be the graph obtained by randomly planting a \kcl in $H$ for some $k=O(1)$, and let $\chi: [n] \to [k]$ be the associated \kcl.
    Then, with high probability, $|\chi^{-1}(a)| = \frac{n}{k} \pm O\Paren{\sqrt{n \log n}}$ for all $a \in [k]$.
\end{lemma}
\begin{proof}
    The proof follows straightforwardly from a Hoeffding bound and a union bound.
\end{proof}

\begin{lemma}[Model entry concentration in random planting]\label{lemma: plantedmodelconcentration}
    Let $H$ be an $n$-vertex $d$-regular graph.
    Let $G$ be the graph obtained by randomly planting a \kcl in $H$ for some $k=O(1)$, and let $\chi: [n] \to [k]$ be the associated \kcl.
    Let $A \in \R^{n \times n}$ be the adjacency matrix of $G$, and let $\tilde{A} = \frac{k}{(k-1)d} A$.
    Let $M(\tilde{A}, \chi)$ be defined as in \cref{def:model}.
    Then, with high probability, $M_{ab}=\frac{1}{k-1} \pm o(1)$ for all $a\neq b\in[k]$.
    % and let $G$ be the graph obtained from $H$ by randomly planting a \kcl in it for some $k=O(1)$. With high probability, the model defined in \cref{def: randomplanted} satisfies $M_{ab}=\frac1{k-1}\pm o(1)$ for any $a\neq b\in[k]$.
\end{lemma}

\begin{proof}
    Let $d' = \frac{(k-1)d}{k}$, and for $x \in [n]$ and $a \in [k]$, let $\D{x}{a} = \sum_{y \in \chi^{-1}(a)} \tilde{A}_{xy}$.
    By linearity of expectation,
    \[
    \begin{aligned}
    \E\Brac{\sum_{x\in \chi^{-1}(a)} \D{x}{b}}
    &= \frac1{d'}\;\E\Brac{\sum_{x\sim y\in E(H)}
        \Paren{\1[\chi(x)=a\land\chi(y)=b]
            +\1[\chi(x)=b\land\chi(y)=a]}} \\[6pt]
    &= \frac1{d'}\;\frac{nd}{2}\;\frac{2}{k^2}
    \;=\;\frac{n}{k(k-1)}\,.
    \end{aligned}
    \]

    By the $d$-regularity of $H$, changing the color of a single vertex can increase or decrease $\sum_{x\in \chi^{-1}(a)} \D{x}{b}$ by at most $O(1)$. Hence, McDiarmid's inequality gives that $\sum_{x\in \chi^{-1}(a)} \D{x}{b}$ is with probability $O(n^{-2})$ within $O(\sqrt{n\log{n}})$ of its expectation. Using \cref{fact: balancedparts} and putting it all together, we get
    $$M_{ab}=\frac{\frac{n}{k(k-1)}\pm O(\sqrt{n\log{n}})}{\frac{n}{k}\pm O(\sqrt{n\log{n}})}=\frac1{k-1}\pm O\Paren{\sqrt{\frac{\log n}{n}}}$$
    with high probability for all $a \neq b \in [k]$.
\end{proof}

Next, we show that with high probability the graph $G$ has a good variance bound.

\begin{lemma}[Variance of random planting]
    \label{lemma: randomplantinglowvariance}
    Let $H$ be an $n$-vertex graph.
    Let $G$ be the graph obtained by randomly planting a \kcl in $H$ for some $k=O(1)$, and let $\chi: [n] \to [k]$ be the associated \kcl.
    Let $A \in \R^{n \times n}$ be the adjacency matrix of $G$, and let $\tilde{A} = \frac{k}{(k-1)d} A$.
    For $x \in [n]$ and $a \in [k]$, let $\D{x}{a} = \sum_{y \in \chi^{-1}(a)} \tilde{A}_{xy}$.
    Then, with high probability, for any $a, b \in [k]$ and $\bm{x}$ uniformly random in $\chi^{-1}(a)$, 
    \[\E_{\bm{x} \sim \chi^{-1}(a)} \Paren{\D{\bm{x}}{b} - \E_{\bm{x} \sim \chi^{-1}(a)} \D{\bm{x}}{b}}^2 \leq 1/d^{\Omega(1)}\,.\]
    % Let $1\leq d < n$ and $k=O(1)$ be parameters. Let $H$ be an $n$-vertex $d$-regular graph and let $G$ be the graph obtained from $H$ by planting a \kcl $\chi$ in it at random.
    % Then, following the definitions of \cref{def: randomplanted}, \whp over $\chi$, for any $a, b \in [k]$ and $\bm{x}$ uniformly random in $\chi^{-1}(a)$,
    % \[\E_{\bm{x} \sim \chi^{-1}(a)} \Paren{\D{\bm{x}}{b} - \E_{\bm{x} \sim \chi^{-1}(a)} \D{\bm{x}}{b}}^2 = o_d(1)\,.\]
\end{lemma}

\begin{proof}
    If $a=b$, we have $\E_{\bm{x} \sim \chi^{-1}(a)} \Paren{\D{\bm{x}}{b} - \E_{\bm{x} \sim \chi^{-1}(a)} \D{\bm{x}}{b}}^2 = 0$ by construction, so we assume $a\neq b$.
    
    Writing $\D{\bm{x}}{b} - \E_{\bm{x} \sim \chi^{-1}(a)} \D{\bm{x}}{b}$ as $\Paren{\D{\bm{x}}{b} - \frac1{k-1}} + \Paren{\frac1{k-1} - \E_{\bm{x} \sim \chi^{-1}(a)} \D{\bm{x}}{b}}$ and using that $(a+b)^2\leq 2(a^2+b^2)$, we get
    \begin{equation}\label{eq:variance_split}
        \E_{\bm{x} \sim \chi^{-1}(a)} \Bigl(\D{\bm{x}}{b} - \E_{\bm{x} \sim \chi^{-1}(a)} \D{\bm{x}}{b}\Bigr)^2
        \;\le\;
        2\Bigl(\E_{\bm{x} \sim \chi^{-1}(a)} \bigl(\D{\bm{x}}{b} - \tfrac1{k-1}\bigr)^2
        +\bigl(\tfrac1{k-1}-M_{ab}\bigr)^2\Bigr).
    \end{equation}

    Let us define $S=\sum_{x \in \chi^{-1}(a)} \Paren{\D{x}{b} - \frac1{k-1}}^2$ to analyze the first term. 
    Let $d' = \frac{(k-1)d}{k}$.
    For any $x\in V(G)$, conditioned on $x\in \chi^{-1}(a)$, by $d$-regularity we have $\D{x}{b}\sim \frac1{d'}\mathrm{Bin}(d, \frac1{k})$.
    Hence, by the law of total probability and linearity of expectation, we get
    $$\E_{\chi}\Brac{S}=\frac{n}{k}\frac1{d'^2}d\frac1{k}\Paren{1-\frac1{k}}=\frac{n}{(k-1)kd} \leq O(n/d)\,.$$ 

    It is easy to verify that changing the color of a single vertex can increase or decrease $S$ by at most $O(1)$. Therefore, by McDiarmid's inequality,
    \[
    \Pr\Brac{\,\bigl|S-\E[S]\bigr|\ge \sqrt{n\log n}}
    \;\le\;
    2\exp\!\Paren{-\Omega\!\Paren{\tfrac{n\log n}{n}}}
    =o(1).
    \]
    Hence, with high probability
    \[
    S = \E[S]\pm O\Paren{\sqrt{n\log n}}
    = O(n/d) \pm O\Paren{\sqrt{n\log n}}.
    \]
    Then, using \cref{fact: balancedparts}, with high probability
    \[\E_{\bm{x} \sim \chi^{-1}(a)} \Paren{\D{\bm{x}}{b} - \frac1{k-1}}^2
    =\frac{S}{|\chi^{-1}(a)|}
    =O\!\Paren{1/d}\;+\;O\!\Paren{\sqrt{\tfrac{\log n}{n}}}
    = 1/d^{\Omega(1)}.
    \]
    
    By \cref{lemma: plantedmodelconcentration}, the second term of \cref{eq:variance_split} is also bounded by $\Paren{\frac1{k-1}-M_{ab}}^2 \leq 1/d^{\Omega(1)}$, finishing the proof.
\end{proof}

The last ingredient to prove \cref{thm: random-planting-partial} is a lemma that shows how low threshold rank is approximately preserved after planting.

\begin{lemma}[Threshold rank kept after planting]\label[lemma]{lemma: threshold_rank_kept}
    % Let $r_1, r_2\geq 0$, $t_1, t_2\geq 1$, $c>0$ and $k\geq 1$ be parameters.
    Let $H$ be an $n$-vertex graph with adjacency matrix $A_H$.
    Let $G$ be the graph obtained by planting a \kcl in $H$ (not necessarily randomly), and let $A_G$ be its adjacency matrix.
    Let $\bar{A}_H:=\frac1{c}A_H$ and $\bar{A}_G:=\frac1{c}A_G$ be adjacency matrices of $H$ and $G$ respectively normalized by the same parameter $c > 0$.
    For $r_1, r_2 \geq 0$ and $t_1, t_2 \geq 1$, assuming that
    \[
    \rank_{\ge r_1}(\bar{A}_H)\le t_1
    \quad\text{and}\quad
    \rank_{\le -r_2}(\bar{A}_H)\le t_2,
    \]
    it holds that
    \[
    \rank_{\le-(r_1+r_2)}(\bar{A}_G)\;\le\;k\,t_1 + t_2,
    \quad
    \rank_{\ge \,r_1+r_2}(\bar{A}_G)\;\le\;t_1 + k\,t_2.
    \]
    \end{lemma}
    
    \begin{proof}
    Decompose \(\bar{A}_G=\bar{A}_H - \bar{A}_F\), where \(\bar{A}_F\) is block‐diagonal over the \(k\) color classes.  
    By Cauchy's interlacing theorem, each block has \(\le t_1\) eigenvalues \(\ge r_1\) and \(\le t_2\) eigenvalues \(\le -r_2\), so \(\rank_{\le -r_1}(-\bar{A}_F)=\rank_{\ge r_1}(\bar{A}_F)\le k\,t_1\) and \(\rank_{\ge r_2}(-\bar{A}_F)=\rank_{\le -r_2}(\bar{A}_F)\le k\,t_2\).
    Hence, applying Weyl's inequalities
    \begin{gather*}
        \rank_{\le -(r_1+r_2)}(\bar{A}_H - \bar{A}_F)\le\rank_{\le -r_2}(\bar{A}_H)+\rank_{\le -r_1}(-\bar{A}_F)\,, \\
        \rank_{\ge r_1+r_2}(\bar{A}_H - \bar{A}_F)\le\rank_{\ge r_1}(\bar{A}_H)+\rank_{\ge r_2}(-\bar{A}_F)\,,
    \end{gather*}
    we get the two claimed bounds.
    \end{proof}
        
We are now ready to prove \cref{thm: random-planting-partial}.

\begin{proof}[Proof of \cref{thm: random-planting-partial}]
    Let $A \in \R^{n \times n}$ be the adjacency matrix of $G$, and let $\tilde{A} = \frac{k}{(k-1)d} A$.
    Also let $d' = \frac{(k-1)d}{k}$.
    Let $M = M(\tilde{A}, \chi)$ be defined as in \cref{def:model}.
    We aim again to apply \cref{thm:part-recovery} with matrices $M$ and $\tilde{A}$, and we verify that its conditions hold.
    
    Let $W = Z M B^{-1} Z^\top$, with the partition matrices $Z, B$ corresponding to the $k$-coloring $\chi$.
    We have trivially that $B^{1/2} M B^{-1/2}$ is symmetric.
    By \cref{lemma: plantedmodelconcentration} and the triangle inequality, $\Norm{M}\leq \Norm{\frac1{k-1}\Paren{\Allones-\Id}}+\Norm{M-\frac1{k-1}\Paren{\Allones-\Id}}=1+o(1)$, where $\Allones \in \R^{k \times k}$ is the matrix with all entries equal to $1$.
    We also have by \cref{fact: balancedparts} that with high probability $|\chi^{-1}(a)| \geq (1/k-o(1))n$ for all $a \in [k]$.
    
    In addition, as the diagonal entries of $M$ are zero and by \cref{lemma: plantedmodelconcentration} the off-diagonal entries are $\frac1{k-1}\pm o(1)$, we have $\Norm{M^{(a)} - M^{(b)}} \geq 2/k$ for any two rows $M^{(a)}$ and $M^{(b)}$ of $M$.
    For the closeness of $\tilde{A}$ and $W$, because of  \cref{fact: balancedparts} and \cref{lemma: randomplantinglowvariance}, we can apply \cref{lem:small-var-to-indic} and get that, for all $a\in [k]$,
    \[\Norm{\tilde{A} \frac{Z_a}{\Norm{Z_a}} - W  \frac{Z_a}{\Norm{Z_a}}}^2 = \frac{1}{d^{\Omega(1)} \Paren{\frac{1}{k}-o(1)}} \leq O_k\Paren{1/d^{\Omega(1)}}\,.\]

    For the low rank of $\tilde{A}$, we have $\rank_{\geq \lambda}(\frac{1}{d}A_H) \leq t$ by assumption, and as $\Norm{\frac1{d}A_H}\leq 1$ because of $d$-regularity, we can apply \cref{cor:spectral_small_to_small} with $\sigma=\lambda$ to get $\rank_{\leq -\sqrt{2\lambda}}(\frac{1}{d}A_H) \leq t/\lambda^2$ as well. 
    Note that, using that $\lambda \leq 1$, we have that $\lambda \leq \sqrt{2 \lambda}$.
    Let $\lambda' = \max\{\sqrt{2\lambda}, \Omega_k(1)\}$ for some small enough $\Omega_k(1)$, such that $\rank_{\geq \lambda'}(\frac{1}{d}A_H)=O_k(t)$ and $\rank_{\leq -\lambda'}(\frac{1}{d}A_H)=O_k(t)$.
    By applying \cref{lemma: threshold_rank_kept} with normalization parameter $d$, we get $\rank_{\geq 2\lambda'}\Paren{\frac1{d}A_G}=O_k(t)$ and $\rank_{\leq -2\lambda'}\Paren{\frac1{d}A_G}=O_k(t)$.
    As $k\geq 3$, we have $\frac{d}{d'}2\lambda'\leq 3\lambda'$, so $\rank_{\geq 3\lambda'}(\tilde{A})+\rank_{\leq -3\lambda'}(\tilde{A})=O_k(t)$.

    % \begin{itemize}
    %     \item By \cref{fact: balancedparts}, $|\chi^{-1}(a)| \geq (1/k-o(1))n$ for all $a \in [k]$.
    %     \item By construction, $B^{1/2} M B^{-1/2}$ is symmetric.
    %     \item As the diagonal entries of $M$ are zero and by \cref{lemma: plantedmodelconcentration} the off-diagonal entries are $\frac1{k-1}\pm o(1)$, we have $\Norm{M^{(a)} - M^{(b)}} \geq 2/k$ for any two rows $M^{(a)}$ and $M^{(b)}$ of $M$.
    %     \item By triangle inequality and \cref{lemma: plantedmodelconcentration}, $\Norm{M}\leq \Norm{\frac1{k-1}\Paren{\Allones-\Id}}+\Norm{M-\frac1{k-1}\Paren{\Allones-\Id}}=1+o(1)$.
    %     \item We have $\rank_{\geq \lambda}(H) \leq t$ by assumption, and as $\Norm{\frac1{d}A_H}\leq 1$ holds by $d$-regularity, we can apply \cref{cor:spectral_small_to_small} with $\sigma=\lambda$ to get $\rank_{\leq -\sqrt{2\lambda}}(H) \leq t/\lambda^2$ as well. Using that $\lambda \leq 1$, we have $\lambda \leq \sqrt{2 \lambda}$, so by $\lambda=\Omega(1)$\andor{ideally remove this assumption} we get $\rank_{\geq \sqrt{2\lambda}}(H)=O(t)$ and $\rank_{\leq -\sqrt{2\lambda}}(H)=O(t)$. As $k=O(1$), by applying \cref{lemma: threshold_rank_kept} with $c=d$, we get $\rank_{\geq 3\sqrt{\lambda}}\Paren{\frac1{d}A_G}=O(t)$ and $\rank_{\leq -3\sqrt{\lambda}}\Paren{\frac1{d}A_G}=O(t)$. As $k\geq 3$, we have $\frac{d}{d'}3\sqrt{\lambda}\leq 5\sqrt{\lambda}$, so $\rank_{\geq 5\sqrt{\lambda}}(\tilde{A})+\rank_{\leq -5\sqrt{\lambda}}(\tilde{A})=O(t)$.
    %     \item Combining \cref{lemma: randomplantinglowvariance} and \cref{fact: balancedparts}, we can use \cref{lem:small-var-to-indic} to get that for all $a\in [k]$,
    %     \[\Norm{\tilde{A} \frac{Z_a}{\Norm{Z_a}} - W  \frac{Z_a}{\Norm{Z_a}}}^2 = o_d(1)/(1/k-o(1)) = o_d(1)\,.\]
    % \end{itemize}

    % Using the assumed upper bound $\lambda\leq \frac1{101k^3}$, we have $5\sqrt{\lambda}\leq (2/k)\sqrt{1/k-o(1)}/4$.
    Then, by the above, with high probability we satisfy the conditions necessary to apply \cref{thm:part-recovery} with $1/k-o(1)$, $1+o(1)$, $2/k$, $O_k(1/d^{\Omega(1)})$, $O(t)$ and $O(\lambda')$ in the roles of $c$, $\zeta$, $\alpha$, $\delta$, $t$ and $\lambda$, respectively.
    Then, we can recover in time $n^{O(1)} \cdot \Paren{O_{k}(1)}^{t}$ a list of $\Paren{O_{k}(1)}^{t}$ $k$-colorings $\hat{\chi}: [n] \to [k]$ such that, for at least one of them, it holds up to a permutation of the color classes that $\mathbb{E}_{\bm{x} \sim \operatorname{Unif}([n])} \1[\hat\chi(x) \neq \chi(x)] \leq O_k(1/d^{\Omega(1)})$.
    Finally, as in the proof of \cref{thm:kalg-partial}, for each candidate $k$-coloring we apply the rounding algorithm in \cref{lemma: generalrounding} on each of its color classes to obtain $k$ independent sets, and return the set of $k$ independent sets that cover most vertices of $G$.
    The rest of the analysis is identical to that in the proof of \cref{thm:kalg-partial}.
\end{proof}

\subsubsection{Rounding to a full coloring}
\label{sec:full-col}

Next, we show that if $H$ actually is a one-sided expander (instead of having just low top threshold rank), and if the degree $d$ is at least a large enough constant, then we can actually obtain a full \kcl by exploiting the full randomness of the model.

We start from the guarantees of \cref{thm: random-planting-partial-list}, which is a version of \cref{thm: random-planting-partial} that omits the last rounding step: instead, it guarantees the recovery of a \emph{list} of candidate $k$-colorings, such that at least one of them is close to the planted $k$-coloring.
We remark that \cref{thm: random-planting-partial} itself does not guarantee that the output $k$-coloring is close to the planted $k$-coloring, while our strategy for obtaining a full coloring relies on this property --- hence our use of \cref{thm: random-planting-partial-list}.
Then, we show how to convert this list of partial colorings to a full coloring.
Our main result is the following:

\begin{theorem}[Result for random planting, full recovery]
    \label{thm: random-planting-full}
    Let $H$ be an $n$-vertex $d$-regular graph with adjacency matrix $A_H$.
    Let $\lambda$ be the second largest eigenvalue of $\frac{1}{d}A_H$.
    % Suppose $\rank_{\geq \lambda}(\frac{1}{d}A_H) \leq t$ for some $0 < \lambda < 1$ and $t \in \mathbb{N}$.
    Let $G$ be the graph obtained by randomly planting a \kcl in $H$, and let $\chi: [n] \to [k]$ be the associated \kcl.
    Then, for $\lambda > 0$ small enough, there exists an algorithm that runs in time $n^{O(1)} \cdot O_{k}(1)$ and outputs with high probability a $(k, 0)$-coloring for $G$.
    % Let $\Omega(1)\leq\lambda\leq \frac1{O(k^3)}$\andor{this $\Omega(1)$ assumption should ideally be removed}, $d\geq O(1)$ and $k=O(1)$ be parameters. Let $H$ be an $n$-vertex $d$-regular $\lambda$-one-sided expander graph, and let $G$ be the graph obtained from $H$ by planting a \kcl in it at random.
    % Then, there exists an algorithm that runs in time $n^{O(1)} \cdot O_{\lambda, k}(1)$ and \whp fully $k$-colors $G$.
\end{theorem}

The proof and algorithm for \cref{thm: random-planting-full} adapt Theorem B.10 of \cite{MR3536556-David16}. For full details, see \cref{app: plantedcoloring}; below we outline the key modifications.


The key difference between our setting and that of \cite{MR3536556-David16} is that they assume two-sided expansion, while we only assume one-sided expansion.
The two ways in which Theorem B.10 of \cite{MR3536556-David16} uses two-sided expansion are via the Expander Mixing Lemma (EML) (used in Lemma B.17, Lemma B.19, and Lemma B.20 of \cite{MR3536556-David16})
$$e(S, T)\leq d|S|\Paren{\frac{|T|}{n}+c\sqrt{\frac{|T|}{|S|}}}$$
and via the vertex expansion lower bound (used in Lemma B.21 of \cite{MR3536556-David16}): for $|S|=\alpha n$, 
$$\Abs{N(S)\setminus S}\geq \Paren{\frac{1}{(1-\alpha)c^2+\alpha}-1}|S|\,.$$

For the former, we note that Lemmas B.19 and B.20 of \cite{MR3536556-David16} use the EML with $S=T$, and that we can avoid Lemma B.17 of \cite{MR3536556-David16} by assuming a large enough (but still constant) $d$. As the lemma below shows, for $S=T$, the EML holds even under one-sided expansion:

\begin{lemma}[Edge density upper bound]
    \label[lemma]{lemma: emlonesided}
    Let $G$ be an $n$-vertex $d$-regular graph with adjacency matrix $A$ satisfying $\lambda_2(\frac{1}{d}A)\leq c$.
    Then, for any set $S \subseteq [n]$,
    $$\Abs{E(G_S)}\leq\frac{d|S|}{2}\Paren{\frac{|S|}{n}+c}\,,$$
    where $E(G_S)$ is the number of edges in the graph induced by $S$.
\end{lemma}

For the latter use of two-sided expansion, we can show under one-sided expansion a bound that is only slightly weaker, and suffices for our purposes by assuming slightly stronger expansion:

\begin{lemma}[Vertex expansion lower bound]
    \label[lemma]{lemma: ssveonesided}
    Let $G$ be an $n$-vertex $d$-regular graph with adjacency matrix $A$ satisfying $\lambda_2(\frac{1}{d}A)\leq c$.
    Then, for any set $S \subseteq [n]$ with $|S|\leq \alpha n$,
    $$\Abs{N(S)\setminus S}\geq \Paren{\frac{1}{c+\sqrt{\alpha}}-1}|S|\,,$$
    where $N(S)$ is the number of vertices outside $S$ adjacent to at least one vertex in $S$.
\end{lemma}

The proofs of \cref{lemma: emlonesided} and \cref{lemma: ssveonesided} are also in \cref{app: plantedcoloring}.

\subsection{Independent sets}

We state now our result for independent sets.

\begin{theorem}[Independent sets, small bottom rank]
\label[theorem]{thm:independentsetrecovery}
    Let $G$ be an $n$-vertex $d$-regular graph with adjacency matrix $A$, and let $\tilde{A} = \frac{1}{d}A$.
    For $0 < \gamma < 1/4$, suppose $G$ contains an independent set $I \subseteq [n]$ of size $\Paren{\frac{1}{2}-\gamma}n$.
    Suppose $\rank_{\leq -\lambda}(\tilde{A}) \leq t$ for some $0 < \lambda < 1$ and $t \in \mathbb{N}$.
    Then there exists an algorithm with runtime $\poly(n) \cdot (\gamma/(1-\lambda))^{-O(t)}$ that returns an independent set of size at least $\Paren{1/2-\gamma - (4+c)\gamma/(1-\lambda)}n$ for any constant $c > 0$.
\end{theorem}
    
By \cref{cor:spectral_small_to_small}, we immediately get a result under the assumption that the top threshold rank is small.

\begin{corollary}[Independent sets, small top rank]
\label[corollary]{cor:independentsetrecovery}
    Let $G$ be an $n$-vertex $d$-regular graph with adjacency matrix $A$, and let $\tilde{A} = \frac{1}{d}A$.
    For $0 < \gamma < 1/4$, suppose $G$ contains an independent set $I \subseteq [n]$ of size $\Paren{\frac{1}{2}-\e}n$.
    Suppose $\rank_{\geq \lambda}(\tilde{A}) \leq t$ for some $0 < \lambda < 1$ and $t \in \mathbb{N}$.
    Then there exists an algorithm with runtime $\poly(n) \cdot (\gamma/(1-\lambda))^{-O(t)}$ that returns an independent set of size at least $\Paren{1/2-\gamma - (4+c)\gamma/(1-\sqrt{\lambda})}n$ for any constant $c > 0$.
\end{corollary}
\begin{proof}
    Because $\rank_{\geq \lambda}(\tilde{A}) \leq t$, we get by applying \cref{cor:spectral_small_to_small} with $\sigma = c$ that $\rank_{\leq -\lambda'}(\tilde{A}) \leq t/c^2$ for any constant $c > 0$, where $\lambda' = \sqrt{\lambda(1-c)+c} \geq \sqrt{\lambda}$.
    The result then follows by \cref{thm:independentsetrecovery}.
\end{proof}

To prove \cref{thm:independentsetrecovery}, we first prove that the indicator vector of the independent set is close to the subspace spanned by the bottom $t$ eigenvectors of the adjacency matrix:

\begin{lemma}[Closeness to bottom eigenvectors]
    \label[lemma]{lemma:strongerprojectionargument}
    In the same setting as \cref{thm:independentsetrecovery}, let $u \in \R^n$ have $u_x = \frac{1}{\sqrt{n}}$ for $x \in I$ and $u_x = -\frac{1}{\sqrt{n}}$ for $x \not\in I$.
    Then, there exists a vector $\tilde{u} \in \R^n$ with $\norm{\tilde{u}} \leq \norm{u}$ and $||\tilde{u}-u||^2\leq \frac{4\gamma}{1-\lambda}$ such that $\tilde{u}$ is in the span of the bottom $t$ eigenvectors of $\tilde{A}$.
\end{lemma}
\begin{proof}
    First, we have 
    \[u^\top \tilde{A} u = \frac{2}{nd} \Paren{ E(\bar{I}, \bar{I}) - E(I, \bar{I}) } = \frac{2}{nd} \Paren{\gamma n d - \Paren{\frac{1}{2}-\gamma}nd } = -(1-4\gamma)\,.\]
    Second, we derive a lower bound on $u^\top \tilde{A} u$.
    Let $P \in \R^{n \times n}$ be the orthogonal projection to the subspace spanned by the bottom $t$ eigenvectors of $\tilde{A}$.
    Then, because $\rank_{\leq -\lambda}(\tilde{A}) \leq t$ and all eigenvalues of $\tilde{A}$ lie in $[-1, 1]$, we get
    \begin{align*}
    u^\top \tilde{A} u
    &\geq -\lambda \cdot \Paren{1-\Norm{Pu}^2} + (Pu)^\top \tilde{A} (Pu)\\
    &\geq -\lambda \cdot \Paren{1-\Norm{Pu}^2} - \Norm{Pu}^2\\
    &= -\lambda - \gamma \Norm{Pu}^2\,.
    \end{align*}
    Then, plugging in $u^\top \tilde{A} u = -(1-4\e)$ and rearranging, we get
    \[\Norm{Pu}^2 \geq \frac{1-\lambda-4\gamma}{1-\lambda} = 1 - \frac{4\gamma}{1-\lambda}\,.\]
    In particular, $\Norm{u - Pu}^2 = 1 - \Norm{Pu}^2 \leq \frac{4\gamma}{1-\lambda}$.
    Letting $\tilde{u} = Pu$ gives the desired conclusion.
\end{proof}

Furthermore, we can round a vector close to the indicator vector of the independent set:
\begin{lemma}[Independent set rounding]
    \label[lemma]{lemma:independentsetrounding}
    Let $G$ be an $n$-vertex graph with vertex set $V$ that contains an independent set $I \subseteq [n]$.
    Let $u \in \R^n$ have $u_x = \frac{1}{\sqrt{n}}$ for $x \in I$ and $u_x = -\frac{1}{\sqrt{n}}$ for $x \notin I$.
    Also let $\hat{u} \in \R^n$ satisfy $\Norm{\hat{u}-u}^2\leq \alpha$ for some $\alpha \geq 0$.
    Then, given as input $G$ and $\hat{u}$, there exists an algorithm with runtime $\poly(n)$ that returns an independent set of size at least $|I|-\alpha n$.
\end{lemma}
\begin{proof}
    Letting $\hat{I} = \Set{x \in V \mid \hat{u}_x \geq 0}$, by $\Norm{\hat{u}-u}^2\leq \alpha$ we have $\Abs{I\triangle \hat{I}}\Paren{\frac{1}{\sqrt{n}}}^2\leq \alpha$.
    Noting that $\hat{I}$ is a superset of the independent set $\hat{I}\cap I$, by applying \cref{lemma: generalrounding} on $\hat{I}$, we get an independent set of size at least $\Abs{\hat{I}\cap I} - \Abs{\hat{I}\setminus \Paren{\hat{I}\cap I}}=\Abs{I} - \Abs{I\triangle \hat{I}}\geq \Abs{I} - \alpha n$.
\end{proof}

Given \cref{lemma:strongerprojectionargument} and \cref{lemma:independentsetrounding}, we prove~\Cref{thm:independentsetrecovery}.

\begin{proof}[Proof of~\Cref{thm:independentsetrecovery}]
We iterate over a net of the subspace spanned by the bottom $t$ eigenvectors of $\bar{A}$ in order to find a vector close to $u$ as defined in~\Cref{lemma:independentsetrounding}, and then use \cref{lemma:independentsetrounding} to round it to an independent set.

By standard arguments, we can construct a $c \sqrt{\gamma/(1-\lambda)}$-cover of the unit ball in the subspace spanned by the bottom $t$ eigenvectors of $\tilde{A}$ in time $(\gamma/(1-\lambda))^{-O(t)}$, for any constant $c > 0$.
The cover will have size $O(\sqrt{\gamma/(1-\lambda)})^{-t}$.
Then, we iterate over all $\hat{u}$ in this cover and for each of them we apply the rounding algorithm in \cref{lemma:independentsetrounding} to construct an independent set.
We return the largest independent set thus obtained.

For correctness, we note by~\Cref{lemma:strongerprojectionargument} that there exists some $\tilde{u}$ in the span of the bottom $t$ eigenvectors with $\norm{\tilde{u}} \leq \norm{u} \leq 1$ that is $2\sqrt{\gamma/(1-\lambda)}$-close to $u$.
Then the $c \sqrt{\gamma/(1-\lambda)}$-net that we construct is guaranteed to find some $\hat{u}$ that is $c \sqrt{\gamma/(1-\lambda)}$-close to $\tilde{u}$ and thus $(2+c)\sqrt{\gamma/(1-\lambda)}$-close to $u$.
Then, by~\Cref{lemma:independentsetrounding}, the returned independent set has size at least $\Paren{1/2-\gamma - (4+c')\gamma/(1-\lambda)}n$. The time complexity is polynomial in $n$ and $(\gamma/(1-\lambda))^{-O(t)}$.
\end{proof}

\subsection{Rounding tools}

First we state a general rounding algorithm that extracts an independent set from an approximate independent set.

\begin{lemma}[Rounding an approximate independent set]
    \label[lemma]{lemma: generalrounding}
    Let $G$ be an $n$-vertex graph and let $I \subseteq [n]$ be an independent set in it.
    Then, given a set $S\supseteq I$, there is a polynomial-time algorithm that returns an independent set $I'$ such that $|I'| \geq |I| - |S \setminus I|$.
\end{lemma}
\begin{proof}
    Note that $S\setminus I$ is a vertex cover of $G_S$.
    Hence, using the $2$-approximation algorithm for the minimum vertex cover problem, we can compute a vertex cover $C$ of $G_S$ in polynomial time such that $|C| \leq 2|S \setminus I|$.
    Then we output $I':=S\setminus C$, which by construction is an independent set with size at least $|S| - 2|S \setminus I|=|I| - |S \setminus I|$, finishing the proof.
\end{proof}

By applying \cref{lemma: generalrounding} to $k$ color classes, we can extract a coloring from an approximate coloring.

\begin{lemma}[Rounding an approximate $k$-coloring]
    \label[lemma]{cor: kcoloringrounding}
    Let $G$ be an $n$-vertex graph and let $\chi: [n] \to [k]$ be a $3$-coloring of its vertices.
    Suppose that $G$ has a vertex cover of all monochromatic edges of size at most $\e n$.
    Then, given $\chi^{-1}(1), \ldots, \chi^{-1}(k)$, there is a polynomial-time algorithm that returns a $(k, 2\e)$-coloring of $G$.
\end{lemma}

\begin{proof}
    We simply run the algorithm from \cref{lemma: generalrounding} on each set $\chi^{-1}(a)$.
    For each $\chi^{-1}(a)$, let $I_a$ be the largest independent set in it.
    Then we obtain an independent set $I'_a$ satisfying $|I'_a| \geq |I_a| - |\chi^{-1}(a) \setminus I_a|$.
    Then 
    \[\sum_{a=1}^k |I'_a| \geq \sum_{a=1}^k |I_a| - \sum_{a=1}^k |\chi^{-1}(a) \setminus I_a| \geq (1 - \e)n - \e n = (1-2\e)n\,.\]
    % Let $S_i$ be the set of vertices assigned color $i$ by $g$ and let $R_i$ be the set of vertices assigned color $i$ by $f$. As $f$ is a proper coloring, $S_i$ is a super set of the independent set $S_i\cap R_i$, so by applying \cref{lemma: generalrounding} to $S_i$ we obtain an independent set $I_i$ such that $|I_i| \geq \Abs{S_i\cap R_i} - \Abs{S_i\setminus \Paren{S_i\cap R_i}}=\Abs{R_i}-\Abs{R_i\triangle S_i}$. Joining the independent sets $I_{[k]}$ and using $\sum_{i=1}^n|R_i|=n$, we obtain a $\Paren{\sum_{i=1}^k |R_i \triangle S_i|}$-almost coloring. As $g$ is a $\gamma$-approximate coloring with respect to $f$, we have that $\sum_{i=1}^k |R_i \triangle S_i| \leq \gamma$, finishing the proof.
\end{proof}
